\part{数列}
\chapter{数列（一）}

\section{数列的初步认识}
\subsection{数列的基础概念}
数列是离散的函数的观念，单调性
\clearpage
\subsection{和与通项}
\clearpage

\section{等差数列}
\subsection{等差数列的通项公式与前\texorpdfstring{$n$}n项和}
累加，迭代
\clearpage

\subsection{等差数列的性质}
最值，证明是等差等比
\clearpage

\section{等比数列}
\subsection{等比数列的通项公式与}
\clearpage

\subsection{等比数列的性质}
\clearpage

\chapter{数列（二）}

\section{数列的通项公式}
\subsection{累加与累乘}
\clearpage
\subsection{待定系数法的辅助数列}
\clearpage
\subsection{不动点与特征根}
\clearpage
\subsection{分奇偶讨论}
\begin{example}
    数列\(\{a_n\}\) 满足 \(a_1 = 2\)，\(a_{n+1} + a_n = 4n - 3\)，求 \(\{a_n\}\) 的通项公式。
\end{example}
\begin{exercise}
    数列\(\{a_n\}\) 满足 \(a_1 = 1\)，\(a_{n+1} + a_n = 2^{n-1}\)，求 \(\{a_n\}\) 的通项公式。
\end{exercise}

奇偶项数列能求通项的条件是两侧下标奇偶性相同。

\begin{example}
数列 \(\{a_n\}\) 满足 \(a_1 = 3\)，\(a_n a_{n+1} = 9 \times 2^{2n-1}\)，求数列 \(\{a_n\}\) 的通项公式。
\end{example}

\begin{example}
    数列 \(\{a_n\}\) 满足 \(a_1=1, a_2 = 2\)，\(a_{n+2} + (-1)^{n+1} a_n = 1 + (-1)^n\)，求数列 \(\{a_n\}\) 的通项公式。
\end{example}

\begin{example}
    \(\{a_n\}\) 满足 \(a_1 = 1\)，\(a_{n+1} = \begin{cases} 
        a_n + 1 & n \text{为奇} \\
        2a_n & n \text{为偶}
        \end{cases}\)，求 \(\{a_n\}\) 的通项公式。
\end{example}
\begin{exercise}
    已知数列 \(\{a_n\}\) 满足 \(a_1 = 1\)，\(a_{n+1} = 
\begin{cases} 
a_n + 1, & n \text{为奇数}, \\
a_n + 2, & n \text{为偶数}.
\end{cases}\)，求数列 \(\{a_n\}\) 的通项公式．

\end{exercise}

\begin{example}
    在数列 \(\{a_n\}\) 中, \(a_1 = 0\)，且对任意 \(k \in \mathbb{N}^*\)，\(a_{2k-1}\)，\(a_{2k}\)，\(a_{2k+1}\) 成等差数列，其公差为 \(d_k\)。

\begin{enumerate}
    \item 若 \(d_k = 2k\)，求 \(\{a_n\}\) 的通项公式并证明 \(a_{2k}\)，\(a_{2k+1}\)，\(a_{2k+2}\) 成等比数列，\(k \in \mathbb{N}^*\)；
    \item 若对任意 \(k \in \mathbb{N}^*\)，\(a_{2k}\)，\(a_{2k+1}\)，\(a_{2k+2}\) 成等比数列，其公比为 \(q_k\)，设 \(q_1 \neq 1\)，证明 \(\left\{\dfrac{1}{q_k - 1}\right\}\) 成等差数列。
\end{enumerate}
\end{example}

\clearpage

\section{数列求和}
\subsection{倒序相加法}
\clearpage
\subsection{错位相减法}
\clearpage
\subsection{裂项相消法}
\clearpage
\subsection{绝对值求和}
\clearpage
\subsection{奇偶性求和}
\begin{example}
    已知\(a_n=(-1)^nn\)，求前\(n\)项和\(S_n\).
\end{example}
\begin{jie}
    当\(n\)为奇数的时候，
    \begin{align*}
        S_n&=-1+2-3+4+\cdots-(n-2)+(n-1)-n\\
           &=(-1+2)+(-3+4)+\cdots+(n-2-(n-1))-n\\
           &=1+1+\cdots+1-n\\
           &=\dfrac{n-1}{2}-n\\
           &=-\dfrac{n+1}{2}
    \end{align*}
    这里为什么是\(\dfrac{n-1}{2}\)个\(1\)呢，这相当于在问\(n\)里面有多少个\(2\)。

    当\(n\)为偶数的时候，\(n-1\)是奇数，故
    \begin{align*}
        S_n&=S_{n-1}+a_n\\
           &=-\dfrac{n}{2}+n\\
           &=\dfrac{n}{2}
    \end{align*}
    \qed
\end{jie}

\begin{exercise}
    已知\(\{a_n\}=(-1)^nn^2\)，求前\(n\)项和\(S_n\).
\end{exercise}
\begin{jie}
    
\end{jie}
\begin{exercise}
    已知\(a_n=(-1)^n 2^{n-1}\)，求前\(n\)项和\(S_n\).
\end{exercise}
\begin{jie}
    \[S_{n}=\left\{\begin{array}{l}
        \frac{2^{n}-1}{3}, n=2 k \\
        \frac{-2^{n}-1}{3}, n=2 k-1
        \end{array}\right.\]
\end{jie}
\begin{example}
    已知\(a_n=(-1)^{n} \dfrac{4 n}{4 n^{2}-1}\)，求前\(n\)项和\(S_n\).
\end{example}

\begin{jie}
    先将\(a_n\)分解\(a_{n}=(-1)^{n} \dfrac{1}{2 n-1}+(-1)^{n} \dfrac{1}{2 n+1}\).
    仍然采取相似的邻项相加求和的思路。

    当\(n\)为偶数的时候，先计算\(a_{n}+a_{n-1}=\dfrac{1}{2 n-1}+\dfrac{1}{2 n+1}+\left(-\dfrac{1}{2 n-3}-\dfrac{1}{2 n-1}\right)=\dfrac{1}{2 n+1}-\dfrac{1}{2 n-3} \)
这符合裂项相消的形式，最初的两项应为\(a_2+a_1=\dfrac{1}{5}-1\)，故\(S_n=\dfrac{1}{2n+1}-1\).

    当\(n\)为奇数的时候，\(n+1\)是偶数，故
    \begin{align*}
        S_n&=S_{n+1}-a_{n+1}\\
            &=\frac{1}{2 n+3}-1-\frac{1}{2 n+1}-\frac{1}{2 n+3}\\
            &=-\frac{1}{2 n+1}-1
    \end{align*}   
    \qed
\end{jie}


\begin{example}
    已知数列 \(\{a_n\}\) 满足，
\[ a_n = 
\begin{cases} 
2n + 1, & n \text{为奇数} \\
3n - 1, & n \text{为偶数}
\end{cases}
\]
求 \(\{a_n\}\) 的前 \(2n\) 项和。
\end{example}
\begin{jie}
    直接考虑\(a_{2n-1}+a_{2n}=2(2n-1)+1+3(2n)-1=10n-2\),其首项是\(a_1+a_2=8\).
    故\(S_n=\dfrac{(8+10n-2)n}{2}=5n^2+3n\)
    \qed
\end{jie}
\begin{exercise}
    \[
    a_{n}=\left\{
        \begin{array}{ll}
        \dfrac{1}{n(n+2)}& n=2 k \\
        2 n-1&n=2 k-1
        \end{array}\right.
\]
求\(S_{n}\)
\end{exercise}
\begin{jie}先注意到\(\dfrac{1}{n(n+2)}=\dfrac{1}{2}\left(\dfrac{1}{n}-\dfrac{1}{n+2}\right)\)

    \(n\)为偶数时，计算\(a_{n}+a_{n-1}=\dfrac{1}{2}\left(\dfrac{1}{n}-\dfrac{1}{n+2}\right)+2 n-3 \),
    其前两项是\(a_1+a_2=\dfrac{1}{2}\left(\dfrac{1}{2}-\dfrac{1}{4}\right)+4-3\).此求和分两部分，前一部分是裂项相消的形式，后一部分是等差数列的形式。

    即\(S_n=\dfrac{1}{2}(1-\dfrac{1}{n+2})+\dfrac{(1+2n-3)\frac{n}{2}}{2}=\dfrac{1}{2}(1-\dfrac{1}{n+2})+\dfrac{n(n-1)}{2}\)

    \(n\)为奇数时，\(S_{n}=S_{n+1}-a_{n+1}=\dfrac{n^{2}+n}{2}+\dfrac{n-1}{4(n+1)}\)

    \qed
\end{jie}
\begin{example}
    数列 \(\{a_n\}\) 中，\(a_n = n^2 \cos \dfrac{2n\pi}{3}\)，求 \(\{a_n\}\) 的前 \(3n\) 项和。
\end{example}
\begin{jie}
    \begin{align*}
        a_{3n-2} & = (3n-2)^2 \cdot \cos\left(\dfrac{(6n-4)\pi}{3}\right) = (3n-2)^2 \cdot \cos\left(-\dfrac{4\pi}{3}\right) = -\dfrac{1}{2}(9n^2 + 4 - 12n) \\
        a_{3n-1} & = (3n-1)^2 \cdot \cos\left(\dfrac{(6n-2)\pi}{3}\right) = (3n-1)^2 \cdot \cos\left(-\dfrac{2\pi}{3}\right) = -\dfrac{1}{2}(9n^2 + 1 - 6n) \\
        a_{3n} & = (3n)^2 \cdot \cos\left(\dfrac{6n\pi}{3}\right) = 9n^2 
    \end{align*}

    故\(a_{3n-2}+a_{3n-1}+a_{3n}=9n - \frac{5}{2}\)

    由此\(S_{3n}=S_{3n}  = \dfrac{\left(9 - \dfrac{5}{2} + 9n - \dfrac{5}{2}\right)n}{2} = \dfrac{(4 + 9n)n}{2}\)
\end{jie}


\begin{example}
    数列 \(\{a_n\}\) 中，\(a_n = 2^n \cos \dfrac{n\pi}{3}\)，求 \(\{a_n\}\) 的前 \(n\) 项和。
\end{example}
\begin{jie}
    由Euler公式，\(\e^{\i\theta}=\cos \theta+\i\sin\theta,\theta\in\mathbb{R}\).
    故\[a_n \cos \frac{n\pi}{3} = 2^n \re\left(\e^{\frac{n\pi}{3}\i}\right) = \re\left((2\e^{\frac{\pi}{3}\i})^n\right) \]
    \begin{align*}
     S_n&= \re\left(\sum_{j=1}^{n} (2\e^{\frac{\pi}{3}\i})^j \right)  \\
        &= \re\left(\frac{(2\e^{\frac{\pi}{3}\i})^{n+1} - 2\e^{\frac{\pi}{3}\i}}{2\e^{\frac{\pi}{3}\i} - 1}\right) \\
        &= \re\left(\frac{2^{n+1} \e^{\frac{(n+1)\pi}{3}\i} - 2\e^{\frac{\pi}{3}\i}}{2\e^{\frac{\pi}{3}\i} - 1}\right) \\
        &= \re\left(\frac{2^{n+1} \cos\left(\frac{(n+1)\pi}{3}\right) + 2^{n+1} \sin\left(\frac{(n+1)\pi}{3}\right)\i - 1}{2\cos\left(\frac{\pi}{3}\right) + 2\sin\left(\frac{\pi}{3}\right)\i - 1} - 1\right) \\
        &= \frac{2^{n+1} \sin\left(\frac{(n+1)\pi}{3}\right)}{2\sin\left(\frac{\pi}{3}\right)} - 1 \\
        &= \frac{\sqrt{3}}{3} \cdot 2^{n+1} \sin\left(\frac{(n+1)\pi}{3}\right) - 1
        \end{align*}

    \qed



\end{jie}




\clearpage

\section{数列的性质}

\subsection{\texorpdfstring{$a_n=f(n)$}s的单调性}
\clearpage
\subsection{\texorpdfstring{$a_{n+1}=f(a_n)$}s的单调性}
\clearpage
\subsection{数列的周期性}
\clearpage

\section{数列放缩初步}
\clearpage
\subsection{拆项放缩}
\clearpage
\subsection{等比放缩}
\clearpage

\section{数学归纳法}
\subsection{Peano公理**}
\clearpage
\subsection{数学归纳法的应用*}
\clearpage